\section{Green Shaping}\label{sec:design}

Green Shaping has three shaping primitives and a controller that chooses their parameters under a joule budget for each node. \Cref{tab:notation} lists the notation used here and in \cref{sec:analysis}.

\begin{table}[t]
\centering
\caption{Notation.}\label{tab:notation}
\footnotesize
\begin{tabular}{@{}ll@{}}
\toprule
Symbol & Meaning \\
\midrule
$N$; $i$ & number of sites; site and link index \\
$\beta$, $r$ & padding base (bytes) and bucket ratio ($r>1$) \\
$W$, $R$ & batching round length (s); reply frame size (bytes) \\
$q_i$ & probability that site $i$ is a cover site in a session \\
$\tau$, $K$, $H$ & DP tick length, ticks per session, $K=\lceil H/\tau\rceil$ \\
$\Delta$ & adjacency bound on bytes entering the DP release \\
$\varepsilon_{\mathrm{dir}}$, $\varepsilon_w$ & budget per link-direction; per tick, $\varepsilon_w=\varepsilon_{\mathrm{dir}}/K$ \\
$a_k$, $Q_k$, $b_k$ & bytes arriving, queued and released at tick $k$ \\
$L_k$ & Laplace noise drawn at tick $k$ \\
$\theta$; $\Theta\supseteq\Theta_1$ & configuration; grid; single-primitive subset \\
$\mathcal{F}(\cdot)$ & configurations within budget and without overflow \\
$\hat A(\theta)$ & attacker advantage on the selection questions \\
$E^{\mathrm{raw}}_n$, $\hat E_n$ & modelled energy; extra energy over unshaped (J) \\
$J_n$ & joule budget of node $n$ per session (J) \\
$e_{\mathrm{tx}}$, $e_{\mathrm{rx}}$ & energy per byte sent and received (J/B) \\
$p_{\mathrm{aw}}$, $t_{\mathrm{tail}}$ & active power above idle (W); activity tail (s) \\
$\pi_i$, $\alpha$ & prior of participation; target detector accuracy \\
$e_{\min,i}$ & minimum extra energy of site $i$ to act as cover in a session \\
\bottomrule
\end{tabular}
\end{table}

\subsection{Placement at the tool-call layer}
The shaper wraps the tool client in the orchestrator and the tool server at each site. At this layer it sees whole calls and replies, knows which session they belong to and knows the full set of sites, information that a packet-level shaper in a middlebox has to infer or does without. It shapes application messages before TLS encryption, so padding and cover messages are encrypted like real traffic and cannot be told apart by content. Tool interfaces such as the Model Context Protocol~\cite{mcp2026spec} define a request and response unit, which gives padding and batching a natural hook. The DP release needs more: it sends a scheduled byte stream, so it runs over one long-lived framed connection per link beneath the tool-call interface (\cref{sec:deploy}). Whether shaping at this layer is cheaper than packet-level shaping at equal attacker advantage is an open empirical question. Our NetShaper comparison measures functional and network costs, not matched privacy.

\subsection{Primitives}

\paragraph{Bucketed padding} Every call and reply body is padded to the next size in a geometric ladder: a message of $m$ bytes becomes
\begin{equation}\label{eq:pad}
\mathrm{pad}(m)=\begin{cases}\beta, & m\le\beta,\\ \beta\, r^{\lceil \log_r (m/\beta)\rceil}, & m>\beta.\end{cases}
\end{equation}
Padding is added inside the encrypted body. The ratio $r$ trades bytes for the number of distinguishable size classes (\cref{prop:pad}). Padding changes neither the number of messages nor their timing.

\paragraph{Synchronised rounds with cover sites} At the start of a session, each site that the agent will not call is made a \emph{cover site} with probability $q_i$; the decision holds for the whole session. Called sites and cover sites form the \emph{active set}. The orchestrator releases calls only at multiples of a round length $W$. At each round boundary every active site receives exactly one call frame, whose size is the round's largest call, padded with \cref{eq:pad}; the orchestrator knows every call it releases, so it can size the frame. Replies travel in frames of a fixed size $R$: every active site sends one frame at the boundary $W$ after the call and, if any reply is incomplete, one more frame per further boundary after a poll from the orchestrator. A called site fills its frames with reply data that is ready; a cover site sends filler. The orchestrator stops polling once every called site has signalled, inside its encrypted frames, that its reply is complete, so all active sites send the same number of frames. Within a session every active link therefore carries the same sequence of times, directions and sizes. What stays visible is the active set, the number of rounds, the call-frame class and the number of reply frames per round, which depends on the largest reply. Rounds cost latency of up to one round per direction and the energy of cover sites, and it removes the per-site timing that otherwise reveals how much work a site performs.

\paragraph{Differentially private release} On each link and direction the shaper keeps a queue and releases bytes only at fixed ticks $k\tau$, $k=1,\dots,K$, for all $N$ sites from the start of the session, independently of the data. At tick $k$ it reads the queue length $Q_k$, draws Laplace noise, and sends exactly
\begin{equation}\label{eq:dp}
b_k=\max\bigl(0,\operatorname{round}(Q_k+L_k)\bigr),\qquad L_k\sim\mathrm{Lap}(\Delta/\varepsilon_w),
\end{equation}
bytes: $\min(Q_k,b_k)$ bytes of queued data followed by padding. $\Delta$ is set on the bytes that enter this primitive, after batching and padding (\cref{prop:order}). This follows the queue-based design of NetShaper~\cite{sabzi2024netshaper}, which noises the amount of data released per fixed interval rather than the send times. NetShaper uses Gaussian noise, adjacency over sliding windows and R\'enyi DP accounting in a middlebox tunnel; we use Laplace noise per link at the end hosts and prove the guarantee for this mechanism directly in \cref{prop:dp} rather than inheriting NetShaper's.

\begin{table*}[t]
\centering
\caption{Where each primitive spends resources and what it protects. Entries are qualitative and follow from the definitions; physical energy coefficients and comparisons across node classes are outside the completed evaluation.}\label{tab:primitives}
\footnotesize
\begin{tabular}{@{}p{2.6cm}p{3.9cm}p{3.6cm}p{2.9cm}p{3.6cm}@{}}
\toprule
Primitive & Extra bytes & Extra awake time & Latency & Guarantee \\
\midrule
Padding & $<(r{-}1)m$ per message if $m>\beta$; $\beta-m$ otherwise & none & none & size classes (\cref{prop:pad}) \\
Rounds with cover sites & frame padding; all traffic of cover sites; polls & cover sites mirror the session; transfers aligned to rounds & up to $W$ per direction and round & equal links within the active set; participation floor (\cref{prop:cover}) \\
DP release & $K\cdot E[\max(0,L_k)]$ per link-direction (\cref{eq:dpcost}) & every tick at every site & queueing, ${\ge}\tau$ & $\varepsilon$-DP (\cref{prop:dp}) \\
\bottomrule
\end{tabular}
\end{table*}

\subsection{Deployment conditions}\label{sec:deploy}
The guarantees of \cref{sec:analysis} hold only if the implementation keeps everything else on the wire independent of the data. (i) \emph{Transport}: each link uses one long-lived TLS connection carrying a framed byte stream (for example an HTTP/2 stream or a WebSocket) on which the tool calls are multiplexed; the DP release writes $b_k$ bytes per tick to this stream. (ii) \emph{Control traffic}: keep-alives, flow-control updates and key updates are disabled during a session or sent on the tick schedule; connection loss aborts the session. (iii) \emph{Session start}: the orchestrator signals the start of a session to all $N$ sites, because a site that starts its ticks only when it is called reveals that it was called; with the DP release every site is therefore active in every session. (iv) \emph{Overflow}: the shaper never sends outside the schedule. A session whose data does not fit in $K$ ticks is truncated and its remaining calls are issued as a new session, which an observer can see; truncations are counted as privacy failures, and $H$ is chosen from calibration sessions to keep them rare. (v) \emph{Concurrency}: sessions on a deployment are serialised; overlapping sessions would need a shared schedule and joint accounting, which we leave open.

\subsection{Energy cost model}
For a shaped session, the controller estimates the energy of node $n$ (a site or the orchestrator) as
\begin{equation}\label{eq:energy}
\begin{aligned}
E^{\mathrm{raw}}_n(\theta)&=e_{\mathrm{tx}}\,B^{\mathrm{tx}}_n+e_{\mathrm{rx}}\,B^{\mathrm{rx}}_n+p_{\mathrm{aw}}\,T^{\mathrm{aw}}_n,\\
\hat E_n(\theta)&=E^{\mathrm{raw}}_n(\theta)-E^{\mathrm{raw}}_n(\varnothing),
\end{aligned}
\end{equation}
where $B^{\mathrm{tx}}_n$ and $B^{\mathrm{rx}}_n$ are the bytes the node sends and receives, $T^{\mathrm{aw}}_n$ is the length of the union of the intervals $[t,t+t_{\mathrm{tail}}]$ over all times $t$ at which the node sends, receives or wakes for a tick, and $\varnothing$ denotes the unshaped session, so $\hat E_n$ is the energy that shaping adds. This is a simple model of a CPU and network interface that stay active for a short time after each transfer. A deployment would fit these four coefficients from instrumented microbenchmarks for its node class. We leave them symbolic: the completed testbed exposes no physical energy counters, so no joule, carbon or measured energy-ranking result is reported (\cref{sec:energy}). The model is a design component, not a fitted empirical predictor in this study. It makes the profiles of the primitives explicit: padding adds only byte terms, rounds add cover traffic but align transfers, and the DP release keeps every node active at every tick.

\subsection{Controller}
\Cref{alg:controller} specifies a proposed energy-budgeted selection procedure, distinct from the functional/cost shortlist used in our implementation study. It selects a configuration offline from sessions recorded on the deployment, split by question into three disjoint sets: attacker training, selection, and a final test set that only the evaluation uses. For each candidate $\theta$ it shapes the sessions, estimates each node's extra energy with \cref{eq:energy}, trains the attacker of \cref{sec:attack} on shaped training sessions and measures its advantage on shaped selection sessions. It returns the feasible configuration with the lowest worst-case advantage over the three secrets. The grid contains every single-primitive configuration, so the selected configuration is never worse on the selection data than the best single primitive within the budget (\cref{rem:dominance}). The controller is rerun when the workload or the budgets change; its cost is $|\Theta|$ shaping passes and attacker trainings.

\begin{algorithm}[t]
\caption{Green Shaping configuration selection}\label{alg:controller}
\footnotesize
\begin{algorithmic}[1]
\Require sessions $\mathcal{C}_{\mathrm{tr}}$, $\mathcal{C}_{\mathrm{sel}}$ with disjoint questions; grid $\Theta\supseteq\Theta_1$; fitted cost models; budgets $(J_n)$
\Ensure configuration $\theta^\star$ or \textsc{infeasible}
\State $\mathcal{F}\gets\emptyset$
\For{$\theta\in\Theta$}
  \State set $\Delta$ of $\theta$ from bytes entering the DP release in $\mathcal{C}_{\mathrm{tr}}\cup\mathcal{C}_{\mathrm{sel}}$
  \State shape $\mathcal{C}_{\mathrm{tr}}$ and $\mathcal{C}_{\mathrm{sel}}$ with $\theta$
  \If{any shaped session overflows} \textbf{continue} \EndIf
  \State $\hat E_n(\theta)\gets$ mean of \cref{eq:energy} over shaped $\mathcal{C}_{\mathrm{sel}}$, for every node $n$
  \If{$\exists n:\hat E_n(\theta)>J_n$} \textbf{continue} \EndIf
  \State train attackers on shaped $\mathcal{C}_{\mathrm{tr}}$
  \State $\hat A(\theta)\gets\max$ over the three secrets of the advantage on shaped $\mathcal{C}_{\mathrm{sel}}$
  \State $\mathcal{F}\gets\mathcal{F}\cup\{\theta\}$
\EndFor
\State \Return $\arg\min_{\theta\in\mathcal{F}}(\hat A(\theta),\max_n \hat E_n(\theta))$ (lexicographic), or \textsc{infeasible} if $\mathcal{F}=\emptyset$
\end{algorithmic}
\end{algorithm}

\subsection{Public-contract admission candidate}\label{sec:admission}
We additionally implement a restricted admission rule using \cref{prop:contract}. Public contracts, risk allocations and policy parameters are frozen before private arrivals. For fixed total privacy budget, we divide $\varepsilon$ among the link-directions and allocate completion risks with total at most $\rho$. A numerical minimax solver distributes these risks to reduce the largest modelled node-cost/budget ratio, using nonnegative node coefficients and the release upper bounds. Every bound and budget is recomputed after solving; failure to admit one candidate is not a proof of global infeasibility. This rule uses analytical privacy and completion certificates rather than the learned attacker objective of Algorithm~\ref{alg:controller}. Expected-cost admission is distinguished from deterministic application-byte admission; neither is a calibrated physical-energy guarantee.

The candidate set must include periodic constant cover. We also implement a public workflow envelope with round cap $R$, planning bound $G$, compute bound $C$ and one-way delivery bound $\ell$. With $T=G+C+2\tau+2\ell$, every site sends a fixed call frame at $u_r=rT+\ell+G+\tau$ and a fixed reply frame at $v_r=(r+1)T$, for $r=0,\ldots,R-1$. Unused sites and rounds still emit cover. Legal bounded calls, computation and planning complete inductively by $RT+\ell$. The abstract frame-metadata transcript is identical across legal private jobs; physical packet timing, TLS behaviour and observable violations require separate validation. Cover must continue after private early completion. We compare this candidate with periodic cover using the same public frame caps, without claiming a new general DP primitive.

\subsection{Reference implementation}
A Python reference implementation provides the three primitives, the cost model, feature extraction and a random-forest attacker, \cref{alg:controller}, and a converter from packet captures to per-link message sequences. Unit tests check \cref{prop:pad,prop:cover,prop:dp,prop:order} numerically, including the exact output distribution of the DP release, a test that fails when its noise is halved, and a test that all active links are identical under rounds with $q_i=1$. The earlier synthetic harness is distinct from the public-data network campaign. The new candidates have synthetic studies in \cref{sec:stress}; matched local TLS acquisition additionally evaluates the workflow envelope and periodic cover in \cref{sec:localtls}. Neither new candidate has been validated on EC2.
